\begin{table}[H]
\centering
\caption{Standardized Effect Sizes for Main Outcomes}
\label{tab:sde}
\begin{adjustbox}{max width=\textwidth}
\begin{tabular}{llcccccl}
\toprule
Outcome & Spec. & $\hat{\beta}$ & SD($X$) & SD($Y$) & SDE & SE(SDE) & Classification \\
\midrule
\multicolumn{8}{l}{\textit{Panel A: Pooled}} \\
In-state tuition & RF & -11.93 & 3.30 & 3,075 & -0.013 & 0.007 & Small negative \\
Enrollment & RF & -320.37 & 3.30 & 13,071 & -0.081 & 0.017 & Moderate negative \\
Completions & RF & -18.50 & 3.30 & 2,486 & -0.025 & 0.008 & Small negative \\
\midrule
\multicolumn{8}{l}{\textit{Panel B: Heterogeneous (by institution type)}} \\
Enrollment (4-year) & RF & -607.49 & 2.97 & 14,107 & -0.128 & 0.022 & Moderate negative \\
Enrollment (2-year) & RF & -133.78 & 3.55 & 11,751 & -0.040 & 0.011 & Small negative \\
\bottomrule
\end{tabular}
\end{adjustbox}
\par\vspace{0.3em}
\begin{minipage}{\textwidth}
{\footnotesize
\textit{Notes:} \textbf{Country:} United States. \textbf{Research question:} Does the \$76 billion Higher Education Emergency Relief Fund (HEERF, 2020--2021), allocated via a Pell-share formula, affect enrollment, tuition, and degree completion at public colleges? \textbf{Policy mechanism:} HEERF allocated emergency funds to U.S. colleges using a formula weighting Pell Grant recipient share and FTE enrollment, with at least 50\% mandated for direct student emergency grants and the remainder available for institutional costs including lost revenue replacement; the formula created quasi-random cross-sectional variation in per-student windfall intensity. \textbf{Outcome definition:} Enrollment is the 12-month unduplicated headcount from IPEDS EFFY; in-state tuition is the annual sticker price for in-state undergraduates from IPEDS IC\_AY; completions are total degrees/certificates awarded from IPEDS C\_A. \textbf{Treatment:} Continuous: predicted HEERF per student in \$1,000s (from the 2018 pre-pandemic Pell share $\times$ FTE allocation formula), interacted with a post-2020 indicator. \textbf{Data:} IPEDS institutional panel, 2015--2022, public 4-year and 2-year institutions, institution-year unit of observation, 12,647 observations. \textbf{Method:} Reduced-form OLS: outcome regressed on predicted HEERF exposure (from pre-pandemic formula) $\times$ post-2020; institution and state $\times$ year fixed effects; standard errors clustered at the institution level. \textbf{Sample:} Public 4-year and 2-year institutions with non-missing tuition and positive enrollment in all years 2015--2022. SDE $= \hat{\beta} \times \text{SD}(X) / \text{SD}(Y)$ where SD($X$) and SD($Y$) are pre-treatment (2015--2019) unconditional standard deviations. Classification refers to magnitude, not statistical significance: Large ($|$SDE$| > 0.15$), Moderate ($0.05$--$0.15$), Small ($0.005$--$0.05$), Null ($< 0.005$).
}
\end{minipage}
\end{table}
